% !TEX root = ../main.tex
%% Corrigés — Chapitre 19 — Procédés sommatoires discrets
%% Résultats vérifiés par verifications/chapitre_19.py.

\corrigesChapitre{19}{Procédés sommatoires discrets}
\corrigesSouples

\begin{blocLivre}{cor-19-1}
\begin{corrige}{1}{Convergence ou divergence ?}
\demarche{Réflexe n°1 : le terme général tend-il vers 0 ?}
\begin{questions}
  \item Géométrique de raison $\dfrac12$ : converge, somme
    $\dfrac{1}{1-1/2}=\res{2}$.
  \item $\dfrac{n}{n+1}\to1\neq0$ : \res{\text{divergence grossière}}.
  \item $(-1)^{n}$ ne tend pas vers $0$ : divergence grossière.
  \item $\dfrac{1}{n(n+1)}=\dfrac1n-\dfrac{1}{n+1}$ : télescopage,
    $S_N=1-\dfrac{1}{N+1}\to\res{1}$.
\end{questions}
\end{corrige}
\end{blocLivre}

\begin{blocLivre}{cor-19-2}
\begin{corrige}{2}{Séries géométriques et exponentielle}
\demarche{Trois sommes du cours.}
$\displaystyle\sum\frac{1}{3^{n}}=\frac{1}{1-1/3}=\res{\frac32}$ ;
$\displaystyle\sum\frac{(-1)^{n}}{2^{n}}=\frac{1}{1+1/2}=\res{\frac23}$ ;
$\displaystyle\sum\frac{1}{n!}=\res{\mathrm{e}}$ — la série exponentielle
en $z=1$.
\end{corrige}
\end{blocLivre}

\begin{blocLivre}{cor-19-3}
\begin{corrige}{3}{Télescopages}
\demarche{Décomposer en différences : tout s'annule sauf les bords.}
\begin{questions}
  \item $\res{1}$ (exercice 1). Et
    $\dfrac{1}{n(n+1)(n+2)}=\dfrac12\Bigl(\dfrac{1}{n(n+1)}
    -\dfrac{1}{(n+1)(n+2)}\Bigr)$ : somme
    $\dfrac12\times\dfrac12=\res{\dfrac14}$.
  \item $\ln\Bigl(1+\dfrac1n\Bigr)=\ln(n+1)-\ln n$ :
    $S_N=\ln(N+1)\to+\infty$ : la série \res{\text{diverge}} — alors que
    son terme général tend vers $0$ : la réciproque de la divergence
    grossière est fausse.
\end{questions}
\end{corrige}
\end{blocLivre}

\begin{blocLivre}{cor-19-4}
\begin{corrige}{4}{Comparer des séries positives}
\demarche{Majorer par une convergente, minorer par une divergente.}
\begin{questions}
  \item Pour $n\geqslant2$ :
    $\dfrac{1}{n^{2}}\leqslant\dfrac{1}{n(n-1)}
    =\dfrac{1}{n-1}-\dfrac1n$, télescopique convergente : les sommes
    partielles de $\sum\frac{1}{n^{2}}$ sont majorées (par $2$) :
    \res{\text{convergence}}.
  \item $\dfrac{1}{n^{2}+n+1}\leqslant\dfrac{1}{n^{2}}$ : converge.
    $\dfrac{1}{\sqrt n}\geqslant\dfrac1n$ et l'harmonique diverge :
    \res{\text{divergence}}.
\end{questions}
\end{corrige}
\end{blocLivre}

\begin{blocLivre}{cor-19-5}
\begin{corrige}{5}{Le critère des équivalents}
\demarche{Un équivalent positif transfère la nature — chapitre 11 au
service des séries.}
\begin{questions}
  \item $\sin\dfrac{1}{n^{2}}\sim\dfrac{1}{n^{2}}$ :
    \res{\text{converge}}.
  \item $1-\cos\dfrac1n\sim\dfrac{1}{2n^{2}}$ : converge.
  \item $\ln\Bigl(1+\dfrac1n\Bigr)\sim\dfrac1n$ : diverge (on retrouve
    l'exercice 3).
  \item $\dfrac{\sqrt{n+1}-\sqrt n}{n}=\dfrac{1}{n(\sqrt{n+1}+\sqrt n)}
    \sim\dfrac{1}{2n^{3/2}}$ : converge (Riemann $\alpha=\frac32$).
\end{questions}
\end{corrige}
\end{blocLivre}

\begin{blocLivre}{cor-19-6}
\begin{corrige}{6}{Les séries de Riemann}
\demarche{Les rectangles encadrent l'intégrale de $t^{-\alpha}$ : la
primitive décide.}
\begin{questions}
  \item Pour $\alpha>1$ :
    $\sum_{n=2}^{N}\dfrac{1}{n^{\alpha}}
    \leqslant\displaystyle\int_{1}^{N}\frac{\mathrm{d}t}{t^{\alpha}}
    \leqslant\dfrac{1}{\alpha-1}$ : sommes partielles majorées,
    convergence. Pour $\alpha\leqslant1$ :
    $\sum_{n=1}^{N}\dfrac{1}{n^{\alpha}}\geqslant
    \displaystyle\int_{1}^{N+1}\frac{\mathrm{d}t}{t^{\alpha}}
    \to+\infty$ : divergence. D'où le critère :
    \res{\text{convergence si et seulement si }\alpha>1}.
  \item $\alpha=2$ et $\alpha=\dfrac32$ : convergentes ;
    $\alpha=\dfrac12$ : divergente.
\end{questions}
\end{corrige}
\end{blocLivre}

\begin{blocLivre}{cor-19-7}
\begin{corrige}{7}{Sommes partielles par série-intégrale}
\demarche{Quand la série diverge, la même comparaison donne l'équivalent
des sommes partielles.}
\begin{questions}
  \item $2\bigl(\sqrt{N+1}-1\bigr)\leqslant
    \displaystyle\sum_{k=1}^{N}\frac{1}{\sqrt k}\leqslant1+2\bigl(\sqrt
    N-1\bigr)$ : en divisant par $2\sqrt N$, les deux bornes tendent vers
    $1$ : $\res{\sum_{k\leqslant N}\tfrac{1}{\sqrt k}\sim2\sqrt N}$.
  \item Le même encadrement avec $\dfrac1t$ donne
    $\ln(N+1)\leqslant H_N\leqslant1+\ln N$ : $H_N\sim\ln N$, et
    $H_N-\ln N$ converge vers $\gamma$ (chapitre 15).
\end{questions}
\end{corrige}
\end{blocLivre}

\begin{blocLivre}{cor-19-8}
\begin{corrige}{8}{Séries alternées : premiers pas}
\demarche{Décroissance vers 0 : le théorème s'applique, et le premier terme
négligé majore l'erreur.}
\begin{questions}
  \item $\dfrac1n$ et $\dfrac{1}{\sqrt n}$ décroissent vers $0$ : les
    deux séries convergent par le théorème des séries alternées.
  \item $|R_N|\leqslant\dfrac{1}{N+1}\leqslant10^{-2}$ dès
    $N\geqslant\res{99}$ ; la somme est $\res{\ln2}$ — vue par les restes
    intégraux aux chapitres 11 et 15, par le théorème général ici.
\end{questions}
\end{corrige}
\end{blocLivre}

\begin{blocLivre}{cor-19-9}
\begin{corrige}{9}{Du lien suite-série}
\demarche{Un produit converge quand la série des logarithmes converge.}
\begin{questions}
  \item $\ln u_n=\displaystyle\sum_{k=1}^{n}\ln\Bigl(1+\frac{1}{k^{2}}\Bigr)$,
    et $\ln\Bigl(1+\dfrac{1}{k^{2}}\Bigr)\sim\dfrac{1}{k^{2}}$ : la série
    converge, donc $(\ln u_n)$ converge et
    $u_n\to\mathrm{e}^{S}>0$ : le produit infini \res{\text{converge}}
    (vers $\dfrac{\operatorname{sh}\pi}{\pi}\approx3{,}676$, pour la
    curiosité).
  \item Tout produit $\prod(1+a_k)$ avec $a_k>-1$ se traite ainsi : le
    logarithme transforme les produits infinis en séries.
\end{questions}
\end{corrige}
\end{blocLivre}

\begin{blocLivre}{cor-19-10}
\begin{corrige}{10}{Restes de séries convergentes}
\demarche{Encadrer le terme entre deux télescopiques voisins.}
\begin{questions}
  \item $\dfrac{1}{n(n+1)}\leqslant\dfrac{1}{n^{2}}\leqslant
    \dfrac{1}{(n-1)n}$ ; en sommant pour $n>N$ :
    $\dfrac{1}{N+1}\leqslant R_N\leqslant\dfrac1N$, d'où
    $\res{R_N\sim\dfrac1N}$ — convergence lente.
  \item $\displaystyle\sum_{n>N}\Bigl(\frac12\Bigr)^{n}
    =\Bigl(\frac12\Bigr)^{N}$ : décroissance \emph{exponentielle},
    incomparablement plus rapide que $\frac1N$ (chapitre 11 :
    croissances comparées).
\end{questions}
\end{corrige}
\end{blocLivre}

\begin{blocLivre}{cor-19-11}
\begin{corrige}{11}{Convergence absolue}
\demarche{La valeur absolue ramène aux séries positives.}
\begin{questions}
  \item $\Bigl|\dfrac{\sin n}{n^{2}}\Bigr|\leqslant\dfrac{1}{n^{2}}$ :
    absolument convergente, donc \res{\text{convergente}} — sans avoir à
    comprendre les signes de $\sin n$.
  \item Non : $\sum\bigl|\tfrac{(-1)^{n-1}}{n}\bigr|=\sum\tfrac1n$
    diverge. Une série convergente non absolument convergente est dite
    \res{\text{semi-convergente}} — espèce fragile, comme le montrera
    l'exercice 15.
\end{questions}
\end{corrige}
\end{blocLivre}

\begin{blocLivre}{cor-19-12}
\begin{corrige}{12}{Le produit de Cauchy et l'exponentielle}
\demarche{Dans le produit de Cauchy des exponentielles, le binôme attend.}
\begin{questions}
  \item $c_n=\displaystyle\sum_{k=0}^{n}\frac{a^{k}}{k!}
    \frac{b^{n-k}}{(n-k)!}
    =\frac{1}{n!}\sum_{k=0}^{n}\binom nk a^{k}b^{n-k}
    =\frac{(a+b)^{n}}{n!}$ : le produit de Cauchy des séries de
    $\mathrm{e}^{a}$ et $\mathrm{e}^{b}$ est la série de
    $\mathrm{e}^{a+b}$, d'où
    $\res{\mathrm{e}^{a}\mathrm{e}^{b}=\mathrm{e}^{a+b}}$.
  \item Ici $c_n=\sum_{k=0}^{n}1=n+1$ :
    $\displaystyle\sum_{n\geqslant0}(n+1)x^{n}
    =\Bigl(\frac{1}{1-x}\Bigr)^{2}$.
\end{questions}
\end{corrige}
\end{blocLivre}

\begin{blocLivre}{cor-19-13}
\begin{corrige}{13}{L'exponentielle complexe, enfin fondée}
\demarche{La convergence absolue est réelle ; le reste est tri de termes.}
\begin{questions}
  \item $\Bigl|\dfrac{z^{n}}{n!}\Bigr|=\dfrac{|z|^{n}}{n!}$, terme d'une
    série convergente : convergence absolue pour tout $z$, et par
    inégalité triangulaire
    $\bigl|\mathrm{e}^{z}\bigr|\leqslant\mathrm{e}^{|z|}$.
  \item Les termes pairs sont réels, les impairs imaginaires purs :
    \[
    \mathrm{e}^{\mathrm{i}\theta}
    =\underbrace{\sum_{p\geqslant0}\frac{(-1)^{p}\theta^{2p}}{(2p)!}}_{\cos\theta}
    +\mathrm{i}\,\underbrace{\sum_{p\geqslant0}
    \frac{(-1)^{p}\theta^{2p+1}}{(2p+1)!}}_{\sin\theta} ,
    \]
    et le morphisme (exercice 12) donne
    $\bigl|\mathrm{e}^{\mathrm{i}\theta}\bigr|^{2}
    =\mathrm{e}^{\mathrm{i}\theta}\mathrm{e}^{-\mathrm{i}\theta}=1$ : le
    cercle trigonométrique du chapitre 3 repose désormais sur des
    fondations complètes.
\end{questions}
\end{corrige}
\end{blocLivre}

\begin{blocLivre}{cor-19-14}
\begin{corrige}{14}{Sommes du type $\sum nx^{n}$}
\demarche{Décaler, soustraire : $(1-x)$ transforme $nx^{n}$ en
géométrique.}
\begin{questions}
  \item $(1-x)\displaystyle\sum_{n\geqslant1}nx^{n}
    =\sum_{n\geqslant1}nx^{n}-\sum_{n\geqslant1}nx^{n+1}
    =\sum_{n\geqslant1}x^{n}=\frac{x}{1-x}$ (les coefficients se
    télescopent), d'où
    $\res{\displaystyle\sum_{n\geqslant1}nx^{n}=\frac{x}{(1-x)^{2}}}$.
  \item En $x=\dfrac12$ : $\res{2}$. Et avec $n^{2}=n(n-1)+n$ :
    $\displaystyle\sum\frac{n^{2}}{2^{n}}
    =\frac14\cdot\frac{2}{(1/2)^{3}}+2=4+2=\res{6}$.
\end{questions}
\end{corrige}
\end{blocLivre}

\begin{blocLivre}{cor-19-15}
\begin{corrige}{15}{Le réarrangement qui change la somme}
\demarche{Les développements de $H_n$ font le calcul ; la semi-convergence
explique le scandale.}
\begin{questions}
  \item La somme des $3N$ premiers termes vaut
    $S_{3N}=\bigl(H_{4N}-\tfrac12H_{2N}\bigr)-\tfrac12H_{N}$ (impairs
    jusqu'à $4N$, pairs jusqu'à $2N$). Avec
    $H_n=\ln n+\gamma+o(1)$ (chapitre 15) :
    \[
    S_{3N}=\ln(4N)-\tfrac12\ln(2N)-\tfrac12\ln N+o(1)
    =2\ln2-\tfrac12\ln2+o(1)\to\res{\tfrac32\ln2} .
    \]
    La même série, réordonnée, a changé de somme.
  \item Une série absolument convergente est une famille sommable : sa
    somme est \emph{invariante par permutation}. La harmonique alternée
    ne l'est pas — la semi-convergence vit de compensations que le
    réarrangement détruit.
\end{questions}
\end{corrige}
\end{blocLivre}

\begin{blocLivre}{cor-19-16}
\begin{corrige}{16}{Fubini positif : premiers calculs}
\demarche{Positif : tous les ordres de sommation sont permis et égaux.}
\begin{questions}
  \item En produit :
    $\Bigl(\sum_i\frac{1}{2^{i}}\Bigr)\Bigl(\sum_j\frac{1}{2^{j}}\Bigr)
    =2\times2=\res{4}$. Par paquets diagonaux $i+j=n$ : le paquet
    contient $n+1$ termes égaux à $\frac{1}{2^{n}}$, d'où
    $\sum_n\dfrac{n+1}{2^{n}}=4$.
  \item La cohérence des deux calculs redonne
    $\displaystyle\sum_{n\geqslant0}\frac{n+1}{2^{n}}=4$ — la formule de
    l'exercice 14, retrouvée sans aucune manipulation de séries.
\end{questions}
\end{corrige}
\end{blocLivre}

\begin{blocLivre}{cor-19-17}
\begin{corrige}{17}{Sommation par paquets}
\demarche{Éléments simples, puis télescopage à pas 2.}
$\dfrac{1}{n^{2}-1}=\dfrac12\Bigl(\dfrac{1}{n-1}-\dfrac{1}{n+1}\Bigr)$
(chapitre 10) : la somme partielle vaut
$\dfrac12\Bigl(1+\dfrac12-\dfrac1N-\dfrac{1}{N+1}\Bigr)\to\res{\dfrac34}$.
Le regroupement des termes distants de $2$ est licite : la série
converge (termes positifs équivalents à $\frac{1}{n^{2}}$), et l'on ne
fait que réécrire ses sommes partielles — ou invoquer la sommation par
paquets d'une famille positive.
\end{corrige}
\end{blocLivre}

\begin{blocLivre}{cor-19-18}
\begin{corrige}{18}{Fubini en action}
\demarche{Un triangle sommé par lignes puis par colonnes.}
\begin{questions}
  \item À $i$ fixé :
    $\displaystyle\sum_{j\geqslant i}x^{j}=\frac{x^{i}}{1-x}$, puis
    $\displaystyle\sum_{i\geqslant0}\frac{x^{i}}{1-x}
    =\frac{1}{(1-x)^{2}}$.
  \item À $j$ fixé, le terme $x^{j}$ apparaît pour
    $i=0,\dots,j$ : contribution $(j+1)x^{j}$, d'où
    $\sum_j(j+1)x^{j}$.
\end{questions}
La famille est positive : Fubini égalise les deux calculs,
$\res{\displaystyle\sum_{j\geqslant0}(j+1)x^{j}=\frac{1}{(1-x)^{2}}}$ —
troisième preuve de la formule, la plus conceptuelle.
\end{corrige}
\end{blocLivre}

\begin{blocLivre}{cor-19-19}
\begin{corrige}{19}{Encadrer les valeurs de $\zeta$}
\demarche{Les rectangles donnent des fourchettes explicites.}
\begin{questions}
  \item $\zeta(2)\leqslant1+\displaystyle\int_{1}^{+\infty}
    \frac{\mathrm{d}t}{t^{2}}=2$ et
    $\zeta(2)\geqslant1+\displaystyle\int_{2}^{+\infty}
    \frac{\mathrm{d}t}{t^{2}}=\frac32$ ; de même
    $\dfrac98\leqslant\zeta(3)\leqslant\dfrac32$.
  \item $\dfrac{\pi^{2}}{6}\approx1{,}645\in\intFF{1{,}5}{2}$ :
    cohérent. (La valeur exacte, le « problème de Bâle », dépasse le
    programme — mais la fourchette, elle, est en notre pouvoir.)
\end{questions}
\end{corrige}
\end{blocLivre}

\begin{blocLivre}{cor-19-20}
\begin{corrige}{20}{Développement asymptotique d'un reste}
\demarche{Des télescopages de plus en plus fins serrent le reste.}
\begin{questions}
  \item Sommer l'encadrement
    $\dfrac{1}{n(n+1)}\leqslant\dfrac{1}{n^{2}}\leqslant\dfrac{1}{(n-1)n}$
    pour $n>N$ : $\dfrac{1}{N+1}\leqslant R_N\leqslant\dfrac1N$, donc
    $NR_N\to1$ : $\res{R_N\sim\dfrac1N}$.
  \item Le télescopage médian donne \emph{exactement}
    $\displaystyle\sum_{n>N}\frac{1}{n^{2}-\frac14}=\frac{1}{N+\frac12}$,
    qui majore $R_N$ : avec a),
    $\dfrac{1}{N+1}<R_N<\dfrac{1}{N+\frac12}$, d'où
    $R_N=\dfrac1N+O\Bigl(\dfrac{1}{N^{2}}\Bigr)$ — un développement
    asymptotique de reste, dans l'esprit du chapitre 11.
\end{questions}
\end{corrige}
\end{blocLivre}

\begin{blocLivre}{cor-19-21}
\begin{corrige}{21}{La constante d'Euler par les séries}
\demarche{La différence entre la somme et l'intégrale, terme à terme.}
\begin{questions}
  \item $\dfrac1k-\ln\Bigl(1+\dfrac1k\Bigr)
    =\dfrac{1}{2k^{2}}+o\Bigl(\dfrac{1}{k^{2}}\Bigr)\geqslant0$
    (chapitre 11) : série convergente par équivalence avec Riemann.
  \item La somme partielle vaut $H_n-\ln(n+1)$, qui tend vers
    $\gamma$ (chapitre 15) : ainsi
    \[
    \res{\gamma=\sum_{k=1}^{+\infty}\Bigl(\frac1k
    -\ln\Bigl(1+\frac1k\Bigr)\Bigr)} :
    \]
    la constante d'Euler, née d'une limite, devient la somme d'une série
    positive — et chaque terme la rapproche par en dessous.
\end{questions}
\end{corrige}
\end{blocLivre}

\begin{blocLivre}{cor-19-22}
\begin{corrige}{22}{Un produit eulérien miniature}
\demarche{L'unicité de la factorisation transforme une somme en produit.}
\begin{questions}
  \item La décomposition en facteurs premiers (chapitre 7) rend
    $(a,b)\mapsto2^{a}3^{b}$ bijective de $\N^{2}$ sur $A$ : par Fubini
    positif,
    \[
    \sum_{n\in A}\frac1n=\sum_{a,b}\frac{1}{2^{a}3^{b}}
    =\Bigl(\sum_a\frac{1}{2^{a}}\Bigr)\Bigl(\sum_b\frac{1}{3^{b}}\Bigr)
    =2\times\frac32=\res{3} :
    \]
    famille sommable, somme exacte.
  \item Avec \emph{tous} les premiers, le même calcul écrirait
    $\sum\frac1n=\prod_p\frac{1}{1-1/p}$ : si les premiers étaient en
    nombre fini, le produit serait fini et l'harmonique convergente —
    absurde. C'est la preuve d'Euler de l'infinité des nombres premiers,
    où l'analyse rejoint l'arithmétique.
\end{questions}
\end{corrige}
\end{blocLivre}

\begin{blocLivre}{cor-19-23}
\begin{corrige}{23}{Le critère de condensation de Cauchy}
\demarche{Dans chaque paquet dyadique, la décroissance encadre tous les
termes par les extrêmes.}
\begin{questions}
  \item Le paquet $\ens{2^{k},\dots,2^{k+1}-1}$ contient $2^{k}$ termes
    compris entre $u_{2^{k+1}}$ et $u_{2^{k}}$ :
    \[
    \tfrac12\cdot2^{k+1}u_{2^{k+1}}\leqslant
    \sum_{n=2^{k}}^{2^{k+1}-1}u_n\leqslant2^{k}u_{2^{k}} .
    \]
    Par sommation par paquets (famille positive), $\sum u_n$ et
    $\sum2^{k}u_{2^{k}}$ ont les mêmes sommes partielles à constantes
    près : \res{\text{même nature}}.
  \item Pour $u_n=n^{-\alpha}$ : la condensée est géométrique de raison
    $2^{1-\alpha}$, d'où le critère de Riemann. Pour
    $u_n=\dfrac{1}{n\ln n}$ : la condensée vaut
    $\sum\dfrac{1}{k\ln2}$, harmonique : $\res{\text{divergence}}$ — une
    série qui diverge plus lentement encore que $\sum\frac1n$.
\end{questions}
\end{corrige}
\end{blocLivre}

\begin{blocLivre}{cor-19-24}
\begin{corrige}{24}{Quand Fubini s'effondre}
\demarche{Un tableau de 1 et de -1 : chaque ordre de sommation raconte une
autre histoire.}
\begin{questions}
  \item Lignes : la ligne $0$ somme à $1$, chaque ligne suivante à
    $1-1=0$ : total $\res{1}$. Colonnes : chaque colonne somme à
    $1-1=0$ : total $\res{0}$. Deux réponses différentes pour « la même
    somme ».
  \item La famille des valeurs absolues contient une infinité de $1$ :
    $\sum|a_{ij}|=+\infty$, la famille n'est \res{\text{pas sommable}} —
    Fubini ne promet rien.
  \item La sommabilité est exactement la frontière : en dessous, l'ordre
    des termes fabrique le résultat (ici $1$ ou $0$, et l'exercice 15
    fabriquait $\frac32\ln2$) ; au-dessus, toutes les routes mènent à la
    même somme. C'est le mot de la fin du chapitre — et la raison d'être
    des familles sommables.
\end{questions}
\end{corrige}
\end{blocLivre}

% !TEX root = ../../main.tex
%% Corrigé du défi du chapitre 19 — vérifié par verifications/defis_16_20.py.

\begin{blocLivre}{cor-19-25a}
\begin{corrige}{25}{Le problème de Bâle : $\zeta(2)=\pi^{2}/6$}
\demarche{Les $m$ nombres $\cot^{2}\frac{r\pi}{2m+1}$ sont les racines d'un
polynôme explicite : Viète donne leur somme ; l'encadrement
$\cot^{2}\theta<\theta^{-2}<1+\cot^{2}\theta$ transmet cette somme exacte à
$\sum1/r^{2}$.}
\partieExo{A}{Une identité polynomiale}
\begin{questions}
  \item Moivre : $(\cos\theta+\mathrm{i}\sin\theta)^{2m+1}
    =\cos(2m+1)\theta+\mathrm{i}\sin(2m+1)\theta$. Par le binôme, la partie
    imaginaire provient des indices impairs $j=2k+1$, avec
    $\mathrm{i}^{2k+1}=(-1)^{k}\mathrm{i}$ : on obtient la formule.
  \item On divise par $\sin^{2m+1}\theta\neq0$ :
    $\frac{\cos^{2m-2k}\theta\,\sin^{2k+1}\theta}{\sin^{2m+1}\theta}
    =\bigl(\cot^{2}\theta\bigr)^{m-k}$, d'où
    $\res{\dfrac{\sin(2m+1)\theta}{\sin^{2m+1}\theta}=P_m\bigl(\cot^{2}\theta\bigr)}$.
  \item Pour $1\leqslant r\leqslant m$, $\theta_r=\frac{r\pi}{2m+1}\in\intOO{0}{\frac\pi2}$
    et $\sin\bigl((2m+1)\theta_r\bigr)=\sin r\pi=0$ : $P_m(x_r)=0$. La
    fonction $\cot^{2}$ est strictement décroissante sur $\intOO{0}{\frac\pi2}$ :
    les $x_r$ sont distincts. $P_m$ est de degré $m$, de coefficient
    dominant $\binom{2m+1}{1}=2m+1$ : ce sont toutes ses racines, et Viète
    donne
    \[
    \sum_{r=1}^{m}x_r=\frac{\binom{2m+1}{3}}{2m+1}=\frac{(2m)(2m-1)}{6}
    =\res{\frac{m(2m-1)}{3}} .
    \]
\end{questions}
\partieExo{B}{L'encadrement}
\begin{questions}[resume]
  \item $\theta-\sin\theta$ et $\tan\theta-\theta$ sont nulles en $0$ et de
    dérivées $1-\cos\theta>0$ et $\tan^{2}\theta>0$ sur
    $\intOO{0}{\frac\pi2}$ : $0<\sin\theta<\theta<\tan\theta$. En passant à
    l'inverse des carrés :
    $\cot^{2}\theta<\frac{1}{\theta^{2}}<\frac{1}{\sin^{2}\theta}=1+\cot^{2}\theta$.
  \item Avec $\frac1{\theta_r^{2}}=\frac{(2m+1)^{2}}{\pi^{2}r^{2}}$ et en
    sommant :
    $\frac{m(2m-1)}{3}<\frac{(2m+1)^{2}}{\pi^{2}}\sum_{r=1}^{m}\frac1{r^{2}}
    <m+\frac{m(2m-1)}{3}=\frac{2m(m+1)}{3}$.
  \item En multipliant par $\frac{\pi^{2}}{(2m+1)^{2}}$, les deux bornes
    $\frac{\pi^{2}m(2m-1)}{3(2m+1)^{2}}$ et
    $\frac{2\pi^{2}m(m+1)}{3(2m+1)^{2}}$ tendent vers
    $\frac{2\pi^{2}}{12}$ : par encadrement,
    $\res{\zeta(2)=\dfrac{\pi^{2}}{6}}$.
\end{questions}
\end{corrige}
\end{blocLivre}

\begin{blocLivre}{cor-19-25b}
\begin{corrigeSuite}
\partieExo{C}{Conséquences}
\begin{questions}[start=7]
  \item Les termes pairs donnent
    $\sum\frac{1}{(2k)^{2}}=\frac{\zeta(2)}{4}$, donc les impairs
    $\frac34\zeta(2)=\res{\frac{\pi^{2}}{8}}$ ; et la série alternée vaut
    $\frac{\pi^{2}}{8}-\frac{\pi^{2}}{24}=\res{\frac{\pi^{2}}{12}}$
    (sommation par paquets de séries absolument convergentes).
  \item $\zeta(2)\in\intOO{S_{10}+\frac1{11}}{S_{10}+\frac{1}{10{,}5}}
    =\intOO{1{,}640\,677}{1{,}645\,006}$, et $\pi=\sqrt{6\zeta(2)}$ :
    \[
    \res{3{,}1375<\pi<3{,}1417},
    \]
    encadrement d'amplitude $0{,}0041$ avec dix termes seulement — grâce au
    contrôle fin du reste. (La vraie valeur $3{,}14159\dots$ est bien dans
    l'intervalle.)
\end{questions}
\end{corrigeSuite}
\end{blocLivre}

